\documentclass[a4paper,11pt]{article}
\usepackage{amsmath,amsthm,amsfonts,amssymb,amscd,amstext,vmargin,graphics,graphicx,tabularx,multicol} 
\usepackage[francais]{babel}
\usepackage[utf8]{inputenc}  
\usepackage[T1]{fontenc} 
\usepackage{pstricks-add,tikz,tkz-tab,variations}
\usepackage[autolanguage,np]{numprint} 

\setmarginsrb{1.5cm}{0.5cm}{1cm}{0.5cm}{0cm}{0cm}{0cm}{0cm} %Gauche, haut, droite, haut
\newcounter{numexo}
\newcommand{\exo}[1]{\stepcounter{numexo}\noindent{\bf Exercice~\thenumexo} : \marginpar{\hfill /#1}}
\reversemarginpar


\newcounter{enumtabi}
\newcounter{enumtaba}
\newcommand{\q}{\stepcounter{enumtabi} \theenumtabi)  }
\newcommand{\qa}{\stepcounter{enumtaba} (\alph{enumtaba}) }
\newcommand{\initq}{\setcounter{enumtabi}{0}}
\newcommand{\initqa}{\setcounter{enumtaba}{0}}

\newcommand{\be}{\begin{enumerate}}
\newcommand{\ee}{\end{enumerate}}
\newcommand{\bi}{\begin{itemize}}
\newcommand{\ei}{\end{itemize}}
\newcommand{\bp}{\begin{pspicture*}}
\newcommand{\ep}{\end{pspicture*}}
\newcommand{\bt}{\begin{tabular}}
\newcommand{\et}{\end{tabular}}
\renewcommand{\tabularxcolumn}[1]{>{\centering}m{#1}} %(colonne m{} centrée, au lieu de p par défault) 
\newcommand{\tnl}{\tabularnewline}

\newcommand{\bmul}[1]{\begin{multicols}{#1}}
\newcommand{\emul}{\end{multicols}}

\newcommand{\trait}{\noindent \rule{\linewidth}{0.2mm}}
\newcommand{\hs}[1]{\hspace{#1}}
\newcommand{\vs}[1]{\vspace{#1}}

\newcommand{\N}{\mathbb{N}}
\newcommand{\Z}{\mathbb{Z}}
\newcommand{\R}{\mathbb{R}}
\newcommand{\C}{\mathbb{C}}
\newcommand{\Dcal}{\mathcal{D}}
\newcommand{\Ccal}{\mathcal{C}}
\newcommand{\mc}{\mathcal}

\newcommand{\vect}[1]{\overrightarrow{#1}}
\newcommand{\ds}{\displaystyle}
\newcommand{\eq}{\quad \Leftrightarrow \quad}
\newcommand{\vecti}{\vec{\imath}}
\newcommand{\vectj}{\vec{\jmath}}
\newcommand{\Oij}{(O;\vec{\imath}, \vec{\jmath})}
\newcommand{\OIJ}{(O;I,J)}


\newcommand{\reponse}[1][1]{%
\multido{}{#1}{\makebox[\linewidth]{\rule[0pt]{0pt}{20pt}\dotfill}
}}

\newcommand{\titre}[5] 
% #1: titre #2: haut gauche #3: bas gauche #4: haut droite #5: bas droite
{
\noindent #2 \hfill #4 \\
#3 \hfill #5

\vspace{-1.6cm}

\begin{center}\rule{6cm}{0.5mm}\end{center}
\vspace{0.2cm}
\begin{center}{\large{\textbf{#1}}}\end{center}
\begin{center}\rule{6cm}{0.5mm}\end{center}
}



\begin{document}
\pagestyle{empty}

\titre{Contrôle - Arithmétique et Scratch}{Nom :}{Prénom :}{Classe}{Date}

\vspace*{0.5cm}

\exo{4}

\q Lequel de ces nombres n'est pas premier ? Justifier votre réponse.\\
19 \hspace*{0.25cm};\hspace*{0.25cm} 23 \hspace*{0.25cm};\hspace*{0.25cm} 93 \hspace*{0.25cm};\hspace*{0.25cm} 43 \hspace*{0.25cm};\hspace*{0.25cm} 61 \hspace*{0.25cm};\hspace*{0.25cm} 89 \hspace*{0.25cm} et \hspace*{0.25cm} 103.\\

\q On donne les nombres suivants :\hfill 36 \hfill ; \hfill 58 180 \hfill; \hfill 27 900 \hfill; \hfill 63 604 \hfill; \hfill 42 324 \hfill; \hfill 34 410.\\
\qa Citer les nombres qui sont divisibles par 5. \textit{(Aucune justification n'est attendue.)}\\
\qa Citer ceux qui sont à la fois divisible par 3 et par 4. \textit{(Aucune justification n'est attendue.)}\\


\q Trouver un nombre à quatre chiffres à la fois divisible par 2;  divisible par 3; divisible par 5 et non divisible par 9. (Aucune justification n'est attendue.)\\



\vspace*{0.25cm}


\exo{3}
\initq \q Citer tous les diviseurs de 48 et 72.\\
\q Quels sont tous \textbf{les diviseurs commun}s à 48 et 72?\\


\vspace*{0.25cm}

\exo{3} 
\initq \q Décomposer 660 et 126 en produit de facteurs premiers. \\
\q En déduire le PGCD de 660 et de 126.\\


\vspace*{0.25cm}



\exo{4}Un pâtissier dispose de 450 morceaux de pommes et de 315 framboises. Il veut préparer le
maximum de tartelettes identiques en utilisant tous les fruits.\\

\noindent \initq \q Peut-il préparer  21 tartelettes ?\\
\q Combien de tartelettes ce pâtissier peut-il faire au maximum ?\\
\q Dans ce cas, combien y aura-t-il de pommes et de framboises sur chacune des tartelettes ?\\

\exo{3}Un carreleur doit poser le carrelage dans une pièce rectangulaire mesurant 6, 48 m de large sur 13, 50 m
de long.
Il souhaite poser des carreaux de carrelage carré et ne faire aucune découpe.\\

$\rightarrow$ Peut-il poser des carreaux de 27 cm de côté ? Justifier votre réponse.\\


\vspace*{0.25cm}

\exo{3} 
\bmul{2}
On considère le script suivant. \\
\initq \q	Combien de segments sont tracés à l’exécution de ce script ?\\
\q	Quelle est la longueur de la ligne tracée ?\\
\q	Lorsque l’on clique sur le drapeau vert, le lutin est au centre de l’écran et est orienté vers le haut. \\
Quelle sera sa position quand l’exécution du script sera terminé ? \textit{Vous donnerez ses coordonnées.}\\


\columnbreak

\includegraphics[scale=1]{scratch.eps} 

\emul





\exo{} BONUS\\
Le nombre de marches d'un escalier est compris entre 40 et 80.\\
- Si on compte ces marches deux par deux, il en reste une.\\
- Si on compte ces marches trois par trois, il en reste deux.\\
- Si on compte ces marches cinq par cinq, il en reste quatre.\\

Quel est le nombre de marches de cet escalier ?\\







\end{document}
